\documentclass[11pt,letterpaper]{article}

% ------------------------------------------------------------
% Page layout and typography
% ------------------------------------------------------------
\usepackage[
    margin=0.72in,
    top=0.65in,
    bottom=0.65in
]{geometry}

\usepackage[T1]{fontenc}
\usepackage[utf8]{inputenc}
\usepackage{lmodern}
\usepackage{microtype}
\usepackage{enumitem}
\usepackage{titlesec}
\usepackage{fancyhdr}
\usepackage{lastpage}
\usepackage{needspace}
\usepackage[hidelinks]{hyperref}

\setlength{\parindent}{0pt}
\setlength{\parskip}{3pt}

\setlist[itemize]{
    leftmargin=1.15em,
    itemsep=2pt,
    topsep=3pt,
    parsep=0pt
}

\setlist[enumerate]{
    leftmargin=2.5em,
    itemsep=6pt,
    topsep=3pt,
    parsep=0pt
}

% Space between section headings and horizontal rules
\titleformat{\section}
    {\large\bfseries}
    {}
    {0pt}
    {}
    [\vspace{4pt}\titlerule]

\titlespacing*{\section}{0pt}{11pt}{8pt}

\titleformat{\subsection}
    {\normalsize\bfseries}
    {}
    {0pt}
    {}

\titlespacing*{\subsection}{0pt}{7pt}{3pt}

% ------------------------------------------------------------
% Footer
% ------------------------------------------------------------
\pagestyle{fancy}
\fancyhf{}
\fancyfoot[L]{\small Durgakant Pushp}
\fancyfoot[R]{\small \thepage\ / \pageref{LastPage}}
\renewcommand{\headrulewidth}{0pt}
\setlength{\footskip}{24pt}

% ------------------------------------------------------------
% Reusable commands
% ------------------------------------------------------------
\newcommand{\me}{\textbf{D. Pushp}}

\newcommand{\entry}[3]{
    \Needspace{4\baselineskip}
    \textbf{#1}\hfill\textbf{#2}\\
    #3\par
}

\newcommand{\award}[2]{
    \textbf{#1}\quad #2\par
}

\begin{document}

% ============================================================
% PAGE 1
% ============================================================

{\LARGE\bfseries Durgakant Pushp}\par

School of Mechanical Engineering, Purdue University
\hfill West Lafayette, Indiana\\
\href{mailto:dpushp@purdue.edu}{dpushp@purdue.edu}
\hfill +1 (812) 974-4282

\section{Research Focus}

Adaptive robotic intelligence through coupled search over
embodiment, representation, and policy under constraints on
information, computation, time, and energy. Research applications
include autonomous racing, intelligent transportation systems,
safe and adaptive autonomous driving, visual navigation, and
heterogeneous robotic systems for field exploration and
infrastructure inspection.

\section{Education}

\entry{Ph.D., Intelligent Systems Engineering}{2021--2026}
{Indiana University Bloomington}

% Confirm the exact title against the deposited dissertation.
Dissertation:
\emph{Bounded Rational Robot: Embodiment, Representation and
the Search for Intelligent Behavior}\\
Advisor: Prof.\ Lantao Liu

\entry{B.Tech., Computer Science and Engineering}{2013--2017}
{West Bengal University of Technology}

\section{Research and Professional Experience}

\entry{Postdoctoral Research Fellow}{2026--Present}
{School of Mechanical Engineering, Purdue University;
Purdue AI Racing (PAIR)}

\begin{itemize}
    \item Mentor a Ph.D.\ student on LiDAR-based localization
    for high-speed autonomous racing; guide research direction
    and technical development of the localization module.

    \item Contributed technical input across perception,
    planning, and control, and helped establish development
    goals and module priorities for PAIR's record-setting
    2026 racing campaign.

    \item Contribute to autonomy-stack architecture and
    integration to support new research capabilities and
    future development.

    \item Extend IJRR research on visual decision-making toward
    controlled stopping during sensor failures in autonomous
    racing; simulation results obtained, with vehicle
    validation planned.

    \item Investigate capability diagnosis and embodiment
    adaptation: how an agent identifies physical limitations
    and evaluates feasible changes under resource constraints.
\end{itemize}

\entry{Team Lead, IU--Luddy Autonomous Racing Team}{2023--2025}
{Indiana University Bloomington}

\begin{itemize}
    \item Led research planning, task allocation, technical
    development, and system integration for a full-scale
    autonomous racing program.

    \item Directed development across localization, perception,
    planning, control, and data analysis for the Dallara AV-24
    platform.

    \item Led and mentored a multidisciplinary team of doctoral,
    master's, undergraduate, and high-school students;
    coordinated development and on-track testing at IMS,
    LVMS, Laguna Seca, and Kentucky Motor Speedway.

    \item Team achieved 128 mph within one month of preparation,
    second place at the 2025 LVMS time trial with a 164 mph
    average, and fourth place at the 2025 Laguna Seca
    road-course event.

    \item Led racing-system development documented in the ISER
    autonomous racing stack publication; contributed sections
    to proposals supporting the program.
\end{itemize}

\newpage

% ============================================================
% PAGE 2
% ============================================================

\section{Research and Professional Experience (continued)}

\entry{Graduate Research Assistant; Technical Project Lead}
{2023--2025}
{Indiana University Bloomington;
Army ERDC-supported robotic inspection of flood-control systems}

\begin{itemize}
    \item Led day-to-day research and project management for
    IU's work on robotic inspection of flood-control
    infrastructure, including research planning, and field experimentation.

    \item Developed a UAV--miniUGV deployable system for
    confined-space inspection and designed and field-tested
    MiniUGV$_2$, integrating mobility, manipulation, sensing,
    and mission autonomy.

    \item Worked within a multi-university collaboration
    involving TAMU, Rice, UT Austin, LSU, MSU, and IU;
    contributed proposal sections for the project.
\end{itemize}

\entry{Graduate Research Assistant; Technical Project Lead}
{2021--2023}
{Indiana University Bloomington;
ARL-supported Off-Road Autonomy / SARA Program}

\begin{itemize}
    \item Led day-to-day technical execution and managed
    research progress for IU's off-road autonomy work,
    including terrain perception and autonomous navigation.

    \item Developed autonomy modules transitioned into ARL's
    Robotics Enterprise software repositories; contributed
    proposal sections supporting the project.
\end{itemize}

\entry{Project Lead, AISUM Autonomous Systems Challenge}
{2021--2022}
{Indiana University Bloomington; U.S.\ Navy CRANE / OUSD(R\&E)}

\begin{itemize}
    \item Developed IU's single-member entry on GPS-denied
    exploration, selected as a national Top-8 finalist;
    prepared proposal sections and challenge reports.
\end{itemize}

\entry{Research Engineer; Field Experiment Lead}{2017--2020}
{TCS Research, Bangalore, India;
collaboration with the Indian Institute of Science}

\begin{itemize}
    \item Developed multi-agent aerial and ground robotic
    systems for industrial automation, including autonomous
    navigation, aerial manipulation, object detection and
    tracking, and decentralized multi-robot coordination.

    \item Conducted research on aerial target detection,
    tracking, and autonomous interception, contributing to
    three peer-reviewed publications at IROS, ICUAS, and ICC.

    \item Led the development, system integration, and outdoor
    field experiments for MBZIRC 2020 Challenge 1, contributing
    to the team's Best Innovation Award.

    \item Contributed to robotic manipulation and adaptive
    gripper development, resulting in a granted Indian patent
    (IN 202321018297, 2024).
\end{itemize}

\section{Research Mentoring and Scholarly Leadership}

\textit{Selected contributions to student research and
publication development.}

\begin{itemize}
    \item \textbf{LiDAR localization for high-speed racing,
    Purdue.}
    Mentor a Ph.D.\ student on the localization module and
    its research development.

    \item \textbf{Seeing Where to Deploy (IROS 2026).}
    Established the research direction and guided the first
    and second authors from project inception through
    manuscript preparation; contributed to study design
    and writing.

    \item \textbf{AFRDA (RA-L 2025).}
    Held weekly research meetings, helped shape the study
    through submission, and contributed to the field experiment and manuscript writing.

    \item \textbf{From Zero to High-Speed Racing (ISER 2025).}
    Led the underlying racing-system development and guided
    research execution and integration; collaborated with
    the first author, who led final manuscript completion.
\end{itemize}

\newpage

% ============================================================
% PAGE 3
% ============================================================

\section{Publications}

\textit{My name is in bold. Preprints are listed separately.}

\subsection{Journal Articles}

\begin{enumerate}[label={[J\arabic*]}]

    \item \me, W. Chen, Z. Chen, C. Luo, J. M. Gregory,
    and L. Liu.
    ``Navigating the Wild: Pareto-Optimal Visual
    Decision-Making in Image Space.''
    \emph{The International Journal of Robotics Research
    (IJRR)}, 2025.\\
    \textbf{Paper of the Season, Q2 2025: Honorable Mention.}

    \item M. A. M. Khan, \me, and L. Liu.
    ``AFRDA: Attentive Feature Refinement for Domain
    Adaptive Semantic Segmentation.''
    \emph{IEEE Robotics and Automation Letters (RA-L)}, 2025.

    \item Z. Chen, \me, J. M. Gregory, and L. Liu.
    ``Pseudo-Trilateral Adversarial Training for Domain
    Adaptive Traversability Prediction.''
    \emph{Autonomous Robots}, vol.\ 47,
    pp.\ 1155--1174, 2023.

\end{enumerate}

\subsection{Conference and Symposium Publications}

\begin{enumerate}[label={[C\arabic*]}]

    \item H. Jardali, I. S. Mohamed, \me, and L. Liu.
    ``Time-Varying Parameter MPC for Robust Path Tracking
    and Lateral Control in Full-Scale Autonomous Racing.''
    \emph{IEEE International Conference on Robotics and
    Automation (ICRA)}, 2026.

    \item S. Lee, S. M. Shithil, \me, L. Liu, and Z. Wang.
    ``Seeing Where to Deploy: Metric RGB-Based Traversability
    Analysis for Aerial-to-Ground Hidden Space Inspection.''
    \emph{IEEE/RSJ International Conference on Intelligent
    Robots and Systems (IROS)}, 2026.

    \item \me, S. Kalhapure, S. M. Shithil, and L. Liu.
    ``MiniUGV$_2$: A Compact UAV-Deployable Tracked Ground
    Vehicle with Manipulation Capabilities.''
    \emph{International Symposium on Experimental Robotics
    (ISER)}, 2025.

    \item H. Jardali, \me, Y. Yu, M. Ali, I. S. Mohamed,
    A. Murillo-Gonzalez, P. D. Coen, M. A.-M. Khan,
    R. C. Pulivendula, S. Park, L. Zhou, and L. Liu.
    ``From Zero to High-Speed Racing: An Autonomous
    Racing Stack.''
    \emph{International Symposium on Experimental Robotics
    (ISER)}, 2025.

    \item \me, J. Xu, Z. Chen, and L. Liu.
    ``Context-Generative Default Policy for Bounded
    Rational Agent.''
    \emph{IEEE/RSJ IROS}, 2024.

    \item M. Ali, \me, Z. Chen, and L. Liu.
    ``Visual-Geometry Gaussian Process-based Navigable
    Space for Autonomous Navigation.''
    \emph{IEEE/RSJ IROS}, 2024.

    \item \me, Z. Chen, C. Luo, J. M. Gregory, and L. Liu.
    ``POVNav: A Pareto-Optimal Mapless Visual Navigator.''
    \emph{ISER}, 2023.

    \item \me, J. Xu, and L. Liu.
    ``Coordination of Bounded Rational Drones through
    Informed Prior Policy.''
    \emph{IEEE/RSJ IROS}, 2023.

    \item \me, S. Kalhapure, K. Das, and L. Liu.
    ``UAV-miniUGV Hybrid System for Hidden Area Exploration
    and Manipulation.''
    \emph{IEEE/RSJ IROS}, pp.\ 1297--1304, 2022.

    \item J. Xu, \me, K. Yin, and L. Liu.
    ``Decision-Making Among Bounded Rational Agents.''
    \emph{International Symposium on Distributed Autonomous
    Robotic Systems (DARS)}, 2022.

    \item Z. Chen, \me, and L. Liu.
    ``CALI: Coarse-to-Fine Alignments for Unsupervised
    Domain Adaptation of Traversability Prediction.''
    \emph{Robotics: Science and Systems (RSS)}, 2022.

    \item U. Rai, \me, A. Choudhary, A. Kumar, B. Chaitanya,
    R. Lima, N. S. Abhinay, R. Majumder, S. Jana, K. Das,
    and D. Ghose.
    ``PARA-FM: Graphical User Interface for Emergency
    Flood Planning and Management.''
    \emph{Proceedings of the 5th World Congress on Disaster
    Management}, vol.\ II, pp.\ 100--106, Routledge, 2022.

\end{enumerate}

\newpage

% ============================================================
% PAGE 4
% ============================================================

\section{Publications (continued)}

\begin{enumerate}[label={[C\arabic*]},start=13]

    \item R. Srivastava, R. Lima, C. Shinde, \me,
    and K. Das.
    ``Estimation and Control for Autonomous UAV System
    to Neutralize Unknown Aerial Maneuvering Targets.''
    \emph{7th Indian Control Conference (ICC)},
    pp.\ 117--122, 2021.

    \item V. Anand, \me, R. Raj, and K. Das.
    ``Gaussian Mixture Model (GMM)-Based Object Detection
    and Tracking Using Dynamic Patch Estimation.''
    \emph{IEEE/RSJ IROS}, pp.\ 4474--4481, 2019.

    \item V. Anand, \me, R. Raj, and K. Das.
    ``Gaussian Mixture Model (GMM)-Based Dynamic Object
    Detection and Tracking.''
    \emph{International Conference on Unmanned Aircraft
    Systems (ICUAS)}, pp.\ 1365--1371, 2019.

\end{enumerate}

\subsection{Preprint}

M. Ali, H. Jardali, Y. Yu, \me, and L. Liu.
``Minimalistic Autonomous Stack for High-Speed Time-Trial
Racing.''
\href{https://arxiv.org/abs/2509.19636}
{arXiv:2509.19636}, 2025.

\subsection{Other Conference Contribution}

\me\ and S. Samaddar.
``A Randomly Guided Algorithm for Exact Pattern Matching.''
National Conference on Recent Advances in Science and Technology
(NCRAST), Heritage Institute of Technology, Kolkata,
September 27--29, 2016.

\section{Honors and Awards}

% Confirm the announcement year of the IJRR recognition.
\award{2026}
{IJRR Paper of the Season Honorable Mention, Q2 2025.}

\award{2024}
{Second Place, Cheng Wu Innovation Challenge;
single-member team.}

\award{2023}
{Audience Choice Award, MassRobotics Form \& Function
Challenge; single-member team.}

\award{2021--2022}
{Top-8 National Finalist, AISUM Challenge;
single-member team.}

\award{2014, 2015, 2017}
{Boeing-sponsored RC aircraft competition awards,
IIT Kharagpur.}

\award{2015, 2016}
{First Prize, RC Aircraft Flying Competition,
NIT Jamshedpur.}

\award{2016}
{Globsyn Faculty Choice Award for analysis of aircraft
crash data.}

\subsection{Team Competition Results}

\award{2026}
{PAIR: First Place, Time Trial; Second Place, Passing
Competition, Indy Autonomous Challenge, Laguna Seca;
record-setting lap.}

\award{2025}
{IU--Luddy: Second Place, LVMS Time Trial; Fourth Place,
Laguna Seca Road Course, Indy Autonomous Challenge.}

\award{2020}
{Best Innovation Award \& \$100K Milestone Grant, Mohamed Bin Zayed International Robotics Challenge}
{Ranked in the top 3 globally during the Phase-1 technical evaluation; served as the lead developer for the core codebase that secured the funding.}

\section{Teaching Experience}

\entry{Guest Lecturer, Autonomous Robots}{August 2024}
{Intelligent Systems Engineering, Indiana University Bloomington}

Undergraduate/graduate course; approximately 20 students.

\begin{itemize}
    \item Delivered a lecture on autonomous racing, covering
    full-stack autonomy, GNSS-based localization, PID and
    Pure Pursuit control, and an introduction to MPPI.
\end{itemize}

\newpage

% ============================================================
% PAGE 5
% ============================================================

\entry{Teaching Assistant, Introduction to Cyber-Physical Systems}
{January--May 2024}
{Intelligent Systems Engineering, Indiana University Bloomington}

Undergraduate course; approximately 30 students.

\begin{itemize}
    \item Supervised student projects using ROS and autonomy
    stacks, held office hours, and supported assessment.
\end{itemize}

\entry{Workshop Organizer and Instructor, Aeromodelling}
{September 2017}
{Jadavpur University, Kolkata, India}

Two-day workshop; approximately 250 undergraduate participants.

\begin{itemize}
    \item Designed and delivered lectures on aerodynamics
    and flight mechanics; guided hands-on construction
    and flight of RC aircraft models.
\end{itemize}

\section{Patent Applications}

\begin{itemize}
    \item R. Lima, \me, et al.
    ``Robot-Mountable Adaptive Gripper Assembly for
    Intercepting Aerial Targets.''
    Indian patent application No.\ 202321018297.
    % If granted, replace with verified grant number and date.

    \item \me, S. Kalhapure, and L. Liu.
    ``UAV Tethered to UGV System.''
    U.S.\ provisional patent application No.\ 63/914,830.
\end{itemize}

\section{Professional Service and Outreach}

\begin{itemize}
    \item Peer reviewer: ICRA, IROS, IEEE Robotics and
    Automation Letters, and ISER.

    \item Invited reviewer, IEEE Intelligent Vehicles
    Symposium (IV), 2026.

    \item Research and robotics mentoring for graduate,
    undergraduate, and high-school students through
    autonomous systems projects.
\end{itemize}

\section{Technical Expertise}

\textbf{Autonomy:}
localization; visual navigation; terrain perception;
planning and control; multi-robot coordination;
field testing and systems integration.\\

\textbf{Programming and robotics:}
C++, Python; ROS 1/2, Gazebo, AirSim, MuJoCo.\\

\textbf{Learning and vision:}
PyTorch, TensorFlow, OpenCV, scikit-learn.\\

\textbf{Sensing:}
LiDAR, radar, RTK-GNSS, IMUs, wheel encoders,
depth and localization cameras, event-based cameras,
sonar, and proximity sensors.

% Submit references separately.
% Purdue Req. 44023 requires at least five references.
% Confirm each referee's current title, affiliation,
% contact information, and willingness to serve.

\end{document}
